\chapter{Smooth Spectral Geometry and the Brave New Motivic Site}
\label{chap:smooth-spectral-motivic-site}

\section{Purpose and standing base convention}
\label{sec:smooth-spectral-purpose}

This chapter applies the cotangent and deformation theory of
\Ref{chap:spectral-deformation-theory} to the geometric classes required
for brave new motivic homotopy theory. It develops differential smoothness
and \`etaleness, their local standard forms and lifting properties, the
smooth finite-presentation site over a quasi-compact quasi-separated
spectral base, the spectral Nisnevich topology, and the brave new affine
line as interval data.

The chapter stops at the geometric input triple
\[
    \left(
        \operatorname{Sm}^{\mathrm{fp}}_{/S},
        \tau_{\mathrm{Nis}},
        \mathbb A^1_S
    \right).
\]
The presheaf localizations imposing Nisnevich descent and
$\mathbb A^1$-invariance belong to the following chapter on brave new
motivic homotopy theory.

All conventions of the preceding chapters remain in force. Unless a
statement explicitly concerns a more general base, the base spectral
scheme $S$ used for the smooth motivic site is assumed quasi-compact and
quasi-separated.

\section{Smooth and \texorpdfstring{\'{e}tale}{etale} spectral morphisms}
\label{sec:smooth-etale-spectral-morphisms}

The geometric morphism classes used in brave new motivic homotopy theory
are defined by finite presentation and the cotangent complex. Smoothness in
this section is differential smoothness. It is not defined by spectral
flatness.

\subsection{Affine definitions}

\begin{definition}[\'{E}tale morphism of connective algebras]
\label{def:etale-spectral-algebra}
A morphism
\[
    A
    \longrightarrow
    B
\]
of connective commutative ring spectra is \textbf{\'{e}tale} if:
\begin{enumerate}[label=(\roman*)]
    \item $B$ is of finite presentation over $A$;

    \item the relative cotangent complex vanishes:
    \[
        L_{B/A}
        \simeq
        0.
    \]
\end{enumerate}
\end{definition}

\begin{definition}[Smooth morphism of connective algebras]
\label{def:smooth-spectral-algebra}
A morphism
\[
    A
    \longrightarrow
    B
\]
of connective commutative ring spectra is \textbf{smooth} if:
\begin{enumerate}[label=(\roman*)]
    \item $B$ is of finite presentation over $A$;

    \item the relative cotangent complex $L_{B/A}$ is a finite locally
    free $B$-module.
\end{enumerate}
\end{definition}

\begin{remark}[Differential and fibre smoothness]
\label{rem:differential-fibre-smoothness}
The smoothness of \Ref{def:smooth-spectral-algebra} is the differential
notion used to define the brave new smooth site. A distinct notion is
obtained by requiring spectral flatness and classical smoothness of
$\pi_0(A)\to\pi_0(B)$. The two notions agree for \`etale morphisms but not
for arbitrary smooth morphisms. In particular, the brave new affine line
is smooth by its cotangent complex even when it is not flat. See
\cite{KhanBraveNewMotivicI}.
\end{remark}

\subsection{Permanence properties}

\begin{proposition}[Base change of smooth and \`etale morphisms]
\label{prop:smooth-etale-base-change}
Let
\[
\begin{tikzcd}[column sep=large, row sep=large]
    A
        \arrow[r]
        \arrow[d]
    &
    B
        \arrow[d]
    \\
    A'
        \arrow[r]
    &
    B':=B\otimes_AA'
\end{tikzcd}
\]
be a pushout square of connective commutative ring spectra.
\begin{enumerate}[label=(\roman*)]
    \item If $A\to B$ is \`etale, then $A'\to B'$ is \`etale.

    \item If $A\to B$ is smooth, then $A'\to B'$ is smooth.
\end{enumerate}
\end{proposition}

\begin{proof}
Finite presentation is preserved under base change by
\Ref{prop:finite-presentation-permanence}. The cotangent base-change
equivalence
\[
    L_{B'/A'}
    \simeq
    L_{B/A}\otimes_BB'
\]
is \Ref{thm:cotangent-base-change}. Vanishing and finite local freeness are
preserved under pullback.
\end{proof}

\begin{proposition}[Composition of smooth and \`etale morphisms]
\label{prop:smooth-etale-composition}
Let
\[
    A
    \longrightarrow
    B
    \longrightarrow
    C
\]
be morphisms of connective commutative ring spectra.
\begin{enumerate}[label=(\roman*)]
    \item If both morphisms are \`etale, then $A\to C$ is \`etale.

    \item If both morphisms are smooth, then $A\to C$ is smooth.
\end{enumerate}
\end{proposition}

\begin{proof}
Finite presentation is stable under composition by
\Ref{prop:finite-presentation-permanence}.

The transitivity triangle is
\[
    L_{B/A}\otimes_BC
    \longrightarrow
    L_{C/A}
    \longrightarrow
    L_{C/B}
    \xrightarrow{\partial}
    (L_{B/A}\otimes_BC)[1].
\]

If both maps are \`etale, the two outer cotangent complexes vanish, and
therefore
\[
    L_{C/A}
    \simeq
    0.
\]
This proves (i).

Suppose both maps are smooth. Put
\[
    P
    :=
    L_{B/A}\otimes_BC,
    \qquad
    Q
    :=
    L_{C/B}.
\]
Both $P$ and $Q$ are finite locally free $C$-modules.

After restricting to a Zariski open cover of $\Spec(C)$, both modules are
finite free. On such an open, the connecting morphism
\[
    \partial
    \colon
    Q
    \longrightarrow
    P[1]
\]
is null. Indeed, if $Q\simeq C^{\oplus r}$, then
\[
    \pi_0
    \Map_{\Mod_C}
    \bigl(
        Q,
        P[1]
    \bigr)
\]
is a finite direct summand of
\[
    \pi_0(P[1])^{\oplus r}
    =
    \pi_{-1}(P)^{\oplus r}
    =
    0,
\]
because $P$ is connective.

The transitivity triangle therefore splits Zariski-locally:
\[
    L_{C/A}
    \simeq
    P\oplus Q.
\]
Thus $L_{C/A}$ is finite locally free. Hence $A\to C$ is smooth.
\end{proof}

\begin{corollary}[\'{E}tale implies smooth]
\label{cor:etale-implies-smooth}
Every \`etale morphism of connective commutative ring spectra is smooth.
\end{corollary}

\begin{proof}
The zero module is finite locally free of rank zero.
\end{proof}

\begin{proposition}[Open immersions are \`etale]
\label{prop:open-immersions-etale}
Every affine spectral open immersion is \`etale.
\end{proposition}

\begin{proof}
Let
\[
    A
    \longrightarrow
    B
\]
represent an affine spectral open immersion. By the definition and
classification of affine spectral open immersions, the morphism is of
finite presentation and is an epimorphism in commutative ring spectra:
the multiplication morphism
\[
    B\otimes_A B
    \longrightarrow
    B
\]
is an equivalence.

Let \(M\) be a \(B\)-module, and let
\[
    B\oplus M
    \longrightarrow
    B
\]
be the associated split square-zero extension. The space of relative
derivations is the homotopy fibre
\[
    \Der_A(B,M)
    \simeq
    \operatorname{fib}_{\id_B}
    \left(
        \Map_{\CAlg(\Sp)_{A/}}
        \bigl(
            B,
            B\oplus M
        \bigr)
        \longrightarrow
        \Map_{\CAlg(\Sp)_{A/}}
        \bigl(
            B,
            B
        \bigr)
    \right).
\]
By extension of scalars along \(A\to B\), this homotopy fibre is
canonically equivalent to
\[
    \operatorname{fib}_{\mu}
    \left(
        \Map_{\CAlg(\Sp)_{B/}}
        \bigl(
            B\otimes_A B,
            B\oplus M
        \bigr)
        \longrightarrow
        \Map_{\CAlg(\Sp)_{B/}}
        \bigl(
            B\otimes_A B,
            B
        \bigr)
    \right),
\]
where
\[
    \mu
    \colon
    B\otimes_A B
    \longrightarrow
    B
\]
is the multiplication morphism. Since \(\mu\) is an equivalence, this
is equivalent to
\[
    \operatorname{fib}_{\id_B}
    \left(
        \Map_{\CAlg(\Sp)_{B/}}
        \bigl(
            B,
            B\oplus M
        \bigr)
        \longrightarrow
        \Map_{\CAlg(\Sp)_{B/}}
        \bigl(
            B,
            B
        \bigr)
    \right).
\]
The object \(B\) is initial in
\[
    \CAlg(\Sp)_{B/},
\]
so both mapping spaces in the last display are contractible. Hence
\[
    \Der_A(B,M)
    \simeq
    *
\]
for every \(B\)-module \(M\).

Since the cotangent complex represents relative derivations, it follows
that
\[
    \Map_{\Mod_B}
    \bigl(
        L_{B/A},
        M
    \bigr)
    \simeq
    *
\]
for every \(B\)-module \(M\). Therefore
\[
    L_{B/A}
    \simeq
    0.
\]
The morphism \(A\to B\) is thus of finite presentation with vanishing
relative cotangent complex, and hence is \`etale. See
\cite{LurieSAG}.
\end{proof}

\subsection{Classical characterization of \texorpdfstring{\'{e}tale}{etale}
morphisms}

\begin{theorem}[Affine spectral \`etale classification]
\label{thm:affine-spectral-etale-classification}
Let
\[
    A
    \in
    \CAlg(\Sp)^{\mathrm{cn}}.
\]
The functor
\[
    B
    \longmapsto
    \pi_0(B)
\]
induces an equivalence from the full subcategory of
\[
    \CAlg(\Sp)^{\mathrm{cn}}_{A/}
\]
spanned by the \`etale \(A\)-algebras to the ordinary category of
\`etale \(\pi_0(A)\)-algebras, regarded as an \(\infty\)-category through
its nerve.

Under this equivalence, the \`etale \(A\)-algebra corresponding to an
ordinary \`etale \(\pi_0(A)\)-algebra \(R\) is flat over \(A\) and
satisfies
\[
    \pi_n(B)
    \simeq
    \pi_n(A)
    \otimes_{\pi_0(A)}
    R
\]
for every integer \(n\).
\end{theorem}

\begin{proof}
This is the affine spectral classification theorem for \`etale
connective commutative ring spectra. We use it as a foundational input
to the spectral \`etale theory developed in this chapter. See
\cite{LurieSAG,KhanBraveNewMotivicI}.
\end{proof}

\begin{theorem}[Spectral \`etale criterion]
\label{thm:spectral-etale-criterion}
For a morphism
\[
    A
    \longrightarrow
    B
\]
of connective commutative ring spectra, the following are equivalent.
\begin{enumerate}[label=(\roman*)]
    \item The morphism \(A\to B\) is \`etale.

    \item The morphism \(A\to B\) is flat and the ordinary ring map
    \[
        \pi_0(A)
        \longrightarrow
        \pi_0(B)
    \]
    is \`etale.
\end{enumerate}
\end{theorem}

\begin{proof}
Suppose first that
\[
    A
    \longrightarrow
    B
\]
is \`etale. By
\Ref{thm:affine-spectral-etale-classification},
the ordinary ring map
\[
    \pi_0(A)
    \longrightarrow
    \pi_0(B)
\]
is \`etale, the \(A\)-algebra \(B\) is flat, and there are natural
isomorphisms
\[
    \pi_n(B)
    \simeq
    \pi_n(A)
    \otimes_{\pi_0(A)}
    \pi_0(B)
\]
for every integer \(n\). This proves
\[
    \textup{(i)}
    \Longrightarrow
    \textup{(ii)}.
\]

Conversely, suppose that \(A\to B\) is flat and that
\[
    \pi_0(A)
    \longrightarrow
    \pi_0(B)
\]
is \`etale. By
\Ref{thm:affine-spectral-etale-classification},
there exists an \`etale connective \(A\)-algebra
\[
    B_{\mathrm{\acute et}}
\]
together with an isomorphism
\[
    \theta
    \colon
    \pi_0(B_{\mathrm{\acute et}})
    \xrightarrow{\simeq}
    \pi_0(B).
\]

The canonical truncation morphism
\[
    B_{\mathrm{\acute et}}
    \longrightarrow
    H\pi_0(B_{\mathrm{\acute et}})
\]
followed by \(H\theta\) gives an \(A\)-algebra morphism
\[
    g_0
    \colon
    B_{\mathrm{\acute et}}
    \longrightarrow
    \tau_{\leq0}B
    \simeq
    H\pi_0(B)
\]
which induces \(\theta\) on degree-zero homotopy.

For every \(n\geq1\), multiplicative Postnikov theory expresses
\[
    \tau_{\leq n}B
    \longrightarrow
    \tau_{\leq n-1}B
\]
as a structured square-zero extension of connective \(A\)-algebras.
Since
\[
    L_{B_{\mathrm{\acute et}}/A}
    \simeq
    0,
\]
the affine cotangent obstruction theory shows that every \(A\)-algebra
morphism
\[
    B_{\mathrm{\acute et}}
    \longrightarrow
    \tau_{\leq n-1}B
\]
admits a lift to
\[
    \tau_{\leq n}B,
\]
and that the space of such lifts is contractible.

Starting with \(g_0\), we therefore obtain a compatible family of
\(A\)-algebra morphisms
\[
    g_n
    \colon
    B_{\mathrm{\acute et}}
    \longrightarrow
    \tau_{\leq n}B.
\]
Postnikov completeness gives
\[
    B
    \simeq
    \lim_{\longleftarrow n}
    \tau_{\leq n}B,
\]
so the compatible family determines an \(A\)-algebra morphism
\[
    g
    \colon
    B_{\mathrm{\acute et}}
    \longrightarrow
    B
\]
whose map on \(\pi_0\) is \(\theta\).

By
\Ref{thm:affine-spectral-etale-classification},
the algebra \(B_{\mathrm{\acute et}}\) is flat over \(A\), and \(B\) is
flat by assumption. Hence, for every integer \(n\), there are natural
identifications
\[
    \pi_n(B_{\mathrm{\acute et}})
    \simeq
    \pi_n(A)
    \otimes_{\pi_0(A)}
    \pi_0(B_{\mathrm{\acute et}})
\]
and
\[
    \pi_n(B)
    \simeq
    \pi_n(A)
    \otimes_{\pi_0(A)}
    \pi_0(B).
\]
Under these identifications, naturality of the flat homotopy-group
formula identifies \(\pi_n(g)\) with
\[
    \id_{\pi_n(A)}
    \otimes
    \pi_0(g).
\]
Since
\[
    \pi_0(g)
    =
    \theta
\]
is an isomorphism, \(\pi_n(g)\) is an isomorphism for every integer
\(n\). Therefore
\[
    g
    \colon
    B_{\mathrm{\acute et}}
    \xrightarrow{\simeq}
    B
\]
is an equivalence.

It follows that \(B\) is an \`etale \(A\)-algebra. Hence
\[
    L_{B/A}
    \simeq
    0
\]
and \(B\) is of finite presentation over \(A\). This proves
\[
    \textup{(ii)}
    \Longrightarrow
    \textup{(i)}.
\]
\end{proof}

\begin{definition}[The \`etale category over a spectral scheme]
\label{def:etale-category-over-spectral-scheme}
For a spectral scheme \(S\), let
\[
    \operatorname{Et}(S)
\]
denote the full subcategory of
\[
    \operatorname{Sch}^{\mathrm{sp}}_{/S}
\]
spanned by the \`etale morphisms
\[
    X
    \longrightarrow
    S.
\]
Its morphisms are arbitrary morphisms of spectral schemes over \(S\).

For an ordinary scheme \(S_0\), let
\[
    \operatorname{Et}(S_0)
\]
denote the ordinary category of ordinary schemes \`etale over \(S_0\),
regarded as an \(\infty\)-category through its nerve. In particular, its
mapping spaces are discrete.
\end{definition}

\begin{theorem}[Invariance of \`etale geometry under classical truncation]
\label{thm:etale-invariance-classical-truncation}
For every spectral scheme \(S\), base change along
\[
    S_{\mathrm{cl}}
    \longrightarrow
    S
\]
induces an equivalence
\[
    \operatorname{Et}(S)
    \xrightarrow{\simeq}
    \operatorname{Et}(S_{\mathrm{cl}}).
\]
This equivalence restricts to the full subcategories spanned by morphisms
of finite presentation.
\end{theorem}

\begin{proof}
The assertion is affine-local on the base. Let
\[
    S
    =
    \Spec(A).
\]
By the spectral classification theorem for \`etale connective
\(A\)-algebras, the functor
\[
    B
    \longmapsto
    \pi_0(B)
\]
induces an equivalence between the \(\infty\)-category of \`etale
connective \(A\)-algebras and the ordinary category of \`etale
\(\pi_0(A)\)-algebras.

Under this equivalence, an \`etale connective \(A\)-algebra \(B\) is
determined by the ordinary \`etale algebra \(\pi_0(B)\), and its higher
homotopy groups are forced by the flat base-change formula
\[
    \pi_n(B)
    \simeq
    \pi_n(A)
    \otimes_{\pi_0(A)}
    \pi_0(B)
\]
for every integer \(n\).

Equivalently, existence and contractible uniqueness may be obtained
inductively along the multiplicative Postnikov tower of \(A\). At each
stage the relevant extension is structured square-zero, and the
vanishing of the relative cotangent complex of an \`etale algebra makes
the lifting space contractible.

Passing to opposite categories gives an equivalence between affine
spectral schemes \`etale over \(\Spec(A)\) and ordinary affine schemes
\`etale over
\[
    \Spec\pi_0(A)
    =
    S_{\mathrm{cl}}.
\]
These affine equivalences are compatible with restriction to affine open
subschemes and with \`etale base change. They therefore glue by \`etale
descent to an equivalence
\[
    \operatorname{Et}(S)
    \xrightarrow{\simeq}
    \operatorname{Et}(S_{\mathrm{cl}}).
\]

Finite presentation is preserved and detected under this equivalence:
affine-locally it is the compactness condition on the corresponding
algebra, and globally the additional quasi-compactness and
quasi-separatedness conditions are unchanged under classical
truncation. Hence the equivalence restricts to the full subcategories of
morphisms of finite presentation.

See \cite{LurieSAG}.
\end{proof}

\begin{proposition}[Diagonals and monomorphisms of \`etale morphisms]
\label{prop:etale-diagonal-and-monomorphisms}
Let
\[
    f
    \colon
    Y
    \longrightarrow
    X
\]
be a morphism of spectral schemes.
\begin{enumerate}[label=(\roman*)]
    \item If \(f\) is \`etale, then its diagonal
    \[
        \Delta_f
        \colon
        Y
        \longrightarrow
        Y\times_XY
    \]
    is an open immersion.

    \item Every \`etale monomorphism of spectral schemes is an open
    immersion.

    \item Every \`etale closed immersion is an open-and-closed
    immersion.
\end{enumerate}
\end{proposition}

\begin{proof}
We first prove (ii). Suppose that
\[
    f
    \colon
    Y
    \longrightarrow
    X
\]
is an \`etale monomorphism. Classical truncation preserves fibre
products, so
\[
    f_{\mathrm{cl}}
    \colon
    Y_{\mathrm{cl}}
    \longrightarrow
    X_{\mathrm{cl}}
\]
is an ordinary \`etale monomorphism. It is therefore an ordinary open
immersion.

Let
\[
    U
    \hookrightarrow
    X
\]
be the spectral open subscheme corresponding to the open subscheme
\(Y_{\mathrm{cl}}\subseteq X_{\mathrm{cl}}\). Then
\[
    U_{\mathrm{cl}}
    \simeq
    Y_{\mathrm{cl}}
\]
over \(X_{\mathrm{cl}}\). Both
\[
    Y
    \longrightarrow
    X
\]
and
\[
    U
    \longrightarrow
    X
\]
are \`etale. By
\Ref{thm:etale-invariance-classical-truncation},
the equivalence of their classical truncations lifts uniquely to an
equivalence
\[
    Y
    \xrightarrow{\simeq}
    U
\]
over \(X\). Thus \(f\) is an open immersion.

We next prove (i). The assertion is local on \(X\) and \(Y\), so consider
an \`etale morphism of connective commutative ring spectra
\[
    A
    \longrightarrow
    B.
\]
Put
\[
    C
    :=
    B\otimes_AB.
\]
The diagonal is represented contravariantly by the multiplication
morphism
\[
    \mu
    \colon
    C
    \longrightarrow
    B.
\]

The morphism
\[
    B
    \longrightarrow
    C
\]
obtained from either tensor factor is a base change of \(A\to B\) and
is therefore \`etale. Consequently the composite
\[
    A
    \longrightarrow
    B
    \longrightarrow
    C
\]
is \`etale. In particular, both \(B\) and \(C\) are compact objects of
\[
    \CAlg(\Sp)_{A/}.
\]

Applying
\Ref{lem:compact-arrow-undercategory}
to
\[
    C
    \xrightarrow{\mu}
    B
\]
in \(\CAlg(\Sp)_{A/}\) shows that \(B\) is compact as a
\(C\)-algebra. Hence \(\mu\) is of finite presentation.

The transitivity triangle for
\[
    A
    \longrightarrow
    C
    \xrightarrow{\mu}
    B
\]
is
\[
    L_{C/A}\otimes_CB
    \longrightarrow
    L_{B/A}
    \longrightarrow
    L_{B/C}.
\]
Both \(A\to C\) and \(A\to B\) are \`etale, so
\[
    L_{C/A}
    \simeq
    0
    \qquad\text{and}\qquad
    L_{B/A}
    \simeq
    0.
\]
It follows that
\[
    L_{B/C}
    \simeq
    0.
\]
Thus
\[
    C
    \longrightarrow
    B
\]
is \`etale.

The diagonal is always a monomorphism. It is therefore an \`etale
monomorphism, and (ii) shows that it is an open immersion. This proves
(i).

Finally, if \(f\) is both \`etale and a closed immersion, then (ii)
shows that it is also an open immersion. Hence it is an open-and-closed
immersion, proving (iii).
\end{proof}

\begin{corollary}[\'{E}tale morphisms are flat]
\label{cor:etale-flat}
Every \`etale morphism of connective commutative ring spectra is flat.
\end{corollary}

\begin{proof}
Apply \Ref{thm:spectral-etale-criterion}.
\end{proof}

\subsection{Brave new affine space and local standard form}

\begin{proposition}[Smoothness of brave new affine space]
\label{prop:brave-new-affine-space-smooth}
For every spectral scheme $X$ and every integer $n\geq0$, the projection
\[
    \mathbb A_X^n
    \longrightarrow
    X
\]
is smooth and of finite presentation.
\end{proposition}

\begin{proof}
The assertion is local on $X$. For $X=\Spec(A)$, the coordinate algebra is
\[
    A\{t_1,\ldots,t_n\}
    =
    \operatorname{Sym}_A(A^{\oplus n}),
\]
which is a finite cell commutative $A$-algebra and hence is of finite
presentation. Its cotangent complex is finite free of rank $n$ by
\Ref{cor:cotangent-brave-new-affine-space}.
\end{proof}

\begin{proposition}[Smoothness of relative vector spaces]
\label{prop:relative-vector-space-smooth}
Let \(X\) be a spectral scheme, and let
\[
    E
    \in
    \QCoh(X)
\]
be finite locally free. Then the projection
\[
    \pi
    \colon
    \mathbb V_X(E)
    :=
    \Spec_X
    \operatorname{Sym}_{\mathcal O_X}(E)
    \longrightarrow
    X
\]
is smooth and of finite presentation. Moreover, there is a canonical
equivalence
\[
    L_{\mathbb V_X(E)/X}
    \simeq
    \pi^*E.
\]
\end{proposition}

\begin{proof}
The assertions are local on \(X\) for the spectral Zariski topology.
Choose an affine open cover
\[
    \{U_\alpha\to X\}_\alpha
\]
on which \(E\) is free of finite rank. Thus, for every \(\alpha\), there
is an integer \(n_\alpha\geq0\) and an equivalence
\[
    E|_{U_\alpha}
    \simeq
    \mathcal O_{U_\alpha}^{\oplus n_\alpha}.
\]

The relative free commutative algebra commutes with restriction, so
\[
\begin{aligned}
    \mathbb V_X(E)\times_XU_\alpha
    &\simeq
    \Spec_{U_\alpha}
    \operatorname{Sym}_{\mathcal O_{U_\alpha}}
    \bigl(
        \mathcal O_{U_\alpha}^{\oplus n_\alpha}
    \bigr)
    \\
    &\simeq
    \mathbb A^{n_\alpha}_{U_\alpha}.
\end{aligned}
\]
By
\Ref{prop:brave-new-affine-space-smooth},
each projection
\[
    \mathbb A^{n_\alpha}_{U_\alpha}
    \longrightarrow
    U_\alpha
\]
is smooth and of finite presentation. Hence
\[
    \mathbb V_X(E)
    \longrightarrow
    X
\]
is smooth and locally of finite presentation.

The projection is affine. It is therefore quasi-compact and
quasi-separated, so it is of finite presentation in the sense of
\Ref{def:finitely-presented-spectral-morphism}.

On each \(U_\alpha\), the cotangent complex of the relative free
commutative algebra gives
\[
\begin{aligned}
    L_{\mathbb V_X(E)/X}
    \big|_{\mathbb V_X(E)\times_XU_\alpha}
    &\simeq
    \mathcal O_{\mathbb A^{n_\alpha}_{U_\alpha}}^{\oplus n_\alpha}
    \\
    &\simeq
    \pi^*
    \bigl(
        E|_{U_\alpha}
    \bigr).
\end{aligned}
\]
These local equivalences are the restrictions of the canonical
free-algebra cotangent equivalence and therefore agree on overlaps.
They descend to a canonical global equivalence
\[
    L_{\mathbb V_X(E)/X}
    \simeq
    \pi^*E.
\]
\end{proof}

\begin{theorem}[Local standard form of smooth morphisms]
\label{thm:smooth-local-standard-form}
Let $A\to B$ be a smooth morphism of connective commutative ring spectra.
For every point of $\Spec\pi_0(B)$, there exists an affine spectral open
neighbourhood $\Spec(B')\to\Spec(B)$, an integer $n\geq0$, and a
factorization
\[
\begin{tikzcd}[column sep=large]
    \Spec(B')
        \arrow[r, "u"]
    &
    \mathbb A_{\Spec(A)}^n
        \arrow[r]
    &
    \Spec(A)
\end{tikzcd}
\]
in which $u$ is \`etale.
\end{theorem}

\begin{proof}
Fix a point of $\Spec\pi_0(B)$. Since $L_{B/A}$ is finite locally free,
after replacing $\Spec(B)$ by an affine open neighbourhood of the chosen
point, we may suppose that
\[
    L_{B/A}
    \simeq
    B^{\oplus n}.
\]
By \Ref{prop:cotangent-connectivity-degree-zero}, there is a natural
identification
\[
    \pi_0(L_{B/A})
    \simeq
    \Omega^1_{\pi_0(B)/\pi_0(A)}.
\]
The ordinary cotangent module is generated by universal differentials.
Choose elements
\[
    b_1,\ldots,b_n
    \in
    \pi_0(B)
\]
whose differentials form a basis after tensoring with the residue field of
the chosen point. The determinant of the resulting morphism
\[
    \pi_0(B)^{\oplus n}
    \longrightarrow
    \pi_0(L_{B/A})
\]
is invertible at that point. Localizing $B$ at this determinant and
relabeling the resulting affine neighbourhood by $B$, we may therefore
suppose that the universal differentials $db_1,\ldots,db_n$ form a basis
of $\pi_0(L_{B/A})$ on the entire neighbourhood.

Put
\[
    P
    :=
    A\{t_1,\ldots,t_n\}
    =
    \operatorname{Sym}_A(A^{\oplus n}).
\]
Choose points
\[
    \widetilde b_i
    \in
    \Omega^\infty B
\]
representing the classes
\[
    b_i
    \in
    \pi_0(B).
\]
They determine an $A$-module morphism
\[
    A^{\oplus n}
    \longrightarrow
    B.
\]
By the free-algebra adjunction, this morphism determines an
$A$-algebra morphism
\[
    \varphi
    \colon
    P
    \longrightarrow
    B.
\]
The transitivity triangle for
\[
    A
    \longrightarrow
    P
    \xrightarrow{\varphi}
    B
\]
is
\[
    L_{P/A}\otimes_PB
    \longrightarrow
    L_{B/A}
    \longrightarrow
    L_{B/P}.
\]
Using
\[
    L_{P/A}
    \simeq
    P^{\oplus n},
\]
the first map is
\[
    B^{\oplus n}
    \longrightarrow
    L_{B/A}
\]
with components $db_i$. By construction it is an isomorphism on
$\pi_0$.

Both its source and target are finite locally free and hence flat
$B$-modules. Therefore, for every integer $m$,
\[
    \pi_m(B^{\oplus n})
    \simeq
    \pi_m(B)
    \otimes_{\pi_0(B)}
    \pi_0(B^{\oplus n})
\]
and
\[
    \pi_m(L_{B/A})
    \simeq
    \pi_m(B)
    \otimes_{\pi_0(B)}
    \pi_0(L_{B/A}).
\]
Under these identifications, the map on $\pi_m$ is obtained by tensoring
the map on $\pi_0$ with $\pi_m(B)$. It is therefore an isomorphism for
every $m$. Hence
\[
    B^{\oplus n}
    \xrightarrow{\simeq}
    L_{B/A},
\]
and the transitivity triangle gives
\[
    L_{B/P}
    \simeq
    0.
\]

It remains to prove finite presentation over $P$. The $A$-algebras $P$
and $B$ are compact in $\CAlg(\Sp)_{A/}$: the first is a finite free cell
algebra, and the second is compact by smoothness. Apply
\Ref{lem:compact-arrow-undercategory} to the morphism
\[
    P
    \longrightarrow
    B
\]
in $\CAlg(\Sp)_{A/}$. It follows that $B$ is compact in
\[
    \bigl(\CAlg(\Sp)_{A/}\bigr)_{P/}
    \simeq
    \CAlg(\Sp)_{P/}.
\]
Thus $P\to B$ is of finite presentation. Together with
$L_{B/P}\simeq0$, this shows that $P\to B$ is \`etale.

Passing to affine spectra gives the required factorization.
\end{proof}

\begin{corollary}[Classical truncation of a smooth morphism]
\label{cor:classical-truncation-smooth}
If $A\to B$ is smooth, then the ordinary ring map
\[
    \pi_0(A)
    \longrightarrow
    \pi_0(B)
\]
is smooth.
\end{corollary}

\begin{proof}
Apply classical truncation to the local standard factorization of
\Ref{thm:smooth-local-standard-form}. The truncation of the \`etale map is
\`etale by \Ref{thm:spectral-etale-criterion}, and the truncation of brave
new affine space is ordinary affine space by
\Ref{prop:classical-truncation-affine-space}.
\end{proof}

\subsection{Global smooth and \texorpdfstring{\'{e}tale}{etale} morphisms}

\begin{definition}[Locally finite presentation]
\label{def:locally-finitely-presented-spectral-morphism}
A morphism
\[
    f
    \colon
    X
    \longrightarrow
    Y
\]
of spectral schemes is \textbf{locally of finite presentation} if, for
every affine open
\[
    V=\Spec(A)
    \longrightarrow
    Y,
\]
there exists an affine open cover
\[
    \{\Spec(B_\alpha)\to X\times_YV\}_{\alpha\in I}
\]
such that every connective commutative $A$-algebra $B_\alpha$ is compact
in
\[
    \CAlg(\Sp)_{A/}.
\]
\end{definition}

\begin{definition}[Finite presentation]
\label{def:finitely-presented-spectral-morphism}
A morphism of spectral schemes is \textbf{of finite presentation} if it
is:
\begin{enumerate}[label=(\roman*)]
    \item locally of finite presentation;

    \item quasi-compact;

    \item quasi-separated.
\end{enumerate}
\end{definition}

\begin{proposition}[Permanence of global finite presentation]
\label{prop:global-finite-presentation-permanence}
Morphisms locally of finite presentation are stable under equivalence,
composition, and base change. Morphisms of finite presentation are stable
under equivalence, composition, and base change.
\end{proposition}

\begin{proof}
Local finite presentation is Zariski-local on both source and target.
After passing to affine charts, its stability under equivalence,
composition, and base change is exactly
\Ref{prop:finite-presentation-permanence}.

Now consider morphisms of finite presentation. Quasi-compact morphisms
and quasi-separated morphisms are each stable under equivalence,
composition, and base change. Combining these permanence properties with
the corresponding permanence of local finite presentation proves that
morphisms of finite presentation are stable under equivalence,
composition, and base change.
\end{proof}

\begin{definition}[Smooth and \`etale morphisms of spectral schemes]
\label{def:global-smooth-etale}
A morphism
\[
    f
    \colon
    X
    \longrightarrow
    Y
\]
of spectral schemes is \textbf{smooth}, respectively \textbf{\`etale}, if
the following affine-local condition holds.

For every affine open
\[
    V=\Spec(A)
    \longrightarrow
    Y
\]
there exists an affine open cover
\[
    \{\Spec(B_\alpha)\to X\times_YV\}_{\alpha\in I}
\]
such that every induced algebra morphism
\[
    A
    \longrightarrow
    B_\alpha
\]
is smooth, respectively \`etale, in the sense of
\Ref{def:smooth-spectral-algebra}, respectively
\Ref{def:etale-spectral-algebra}.
\end{definition}

\begin{proposition}[Locality of smoothness and \`etaleness]
\label{prop:smooth-etale-locality}
The definition of \Ref{def:global-smooth-etale} is independent of the
chosen affine covers. Smoothness and \`etaleness are local on both source
and target for the spectral Zariski topology.
\end{proposition}

\begin{proof}
Finite presentation and finite local freeness of the cotangent complex are
Zariski-local. The cotangent complex itself is compatible with affine
restriction by \Ref{thm:global-cotangent-affine-compatibility}. Vanishing is likewise
local because affine restriction functors over a cover are jointly
conservative.
\end{proof}

\begin{proposition}[Global permanence]
\label{prop:global-smooth-etale-permanence}
Smooth and \`etale morphisms of spectral schemes are stable under
equivalence, composition, and base change. Open immersions are \`etale,
and \`etale morphisms are smooth and flat.
\end{proposition}

\begin{proof}
All assertions are local on source and target. Reduce to the affine
permanence results
\Ref{prop:smooth-etale-base-change},
\Ref{prop:smooth-etale-composition},
\Ref{prop:open-immersions-etale}, and
\Ref{cor:etale-flat}.
\end{proof}



\section{Lifting and neighbourhood theorems}
\label{sec:smooth-etale-lifting-neighbourhoods}

\subsection{Lifting over square-zero extensions}

\begin{corollary}[Lifting \`etale algebras]
\label{cor:lifting-etale-algebras}
In the situation of \Ref{thm:affine-cotangent-obstruction}, suppose that
\[
    A
    \longrightarrow
    B
\]
is \`etale. Then the \(\infty\)-groupoid of pairs
\[
    (B',\alpha),
\]
where \(B'\) is an \`etale \(A'\)-algebra and
\[
    \alpha
    \colon
    B'\otimes_{A'}A
    \xrightarrow{\simeq}
    B
\]
is an equivalence of \(A\)-algebras, is contractible.
\end{corollary}

\begin{proof}
The morphism
\[
    \pi_0(A')
    \longrightarrow
    \pi_0(A)
\]
is a square-zero extension of ordinary commutative rings. Since
\[
    \pi_0(A)
    \longrightarrow
    \pi_0(B)
\]
is \`etale, invariance of ordinary \`etale algebras under nilpotent
extensions supplies an \`etale \(\pi_0(A')\)-algebra \(B'_0\), unique up
to unique isomorphism, together with an isomorphism
\[
    \beta_0
    \colon
    B'_0
    \otimes_{\pi_0(A')}
    \pi_0(A)
    \xrightarrow{\simeq}
    \pi_0(B).
\]

By
\Ref{thm:etale-invariance-classical-truncation},
the algebra \(B'_0\) determines an \`etale connective \(A'\)-algebra
\[
    A'
    \longrightarrow
    B',
\]
unique up to a contractible space of choices.

The base change
\[
    B'\otimes_{A'}A
\]
is an \`etale \(A\)-algebra. Its degree-zero homotopy ring is naturally
identified with
\[
    B'_0
    \otimes_{\pi_0(A')}
    \pi_0(A),
\]
and hence \(\beta_0\) gives an isomorphism
\[
    \pi_0
    \bigl(
        B'\otimes_{A'}A
    \bigr)
    \xrightarrow{\simeq}
    \pi_0(B).
\]

Applying
\Ref{thm:etale-invariance-classical-truncation}
over \(A\), this isomorphism lifts to an equivalence of \(A\)-algebras
\[
    \alpha
    \colon
    B'\otimes_{A'}A
    \xrightarrow{\simeq}
    B.
\]
The space of equivalences inducing the specified isomorphism on
degree-zero homotopy is contractible.

Ordinary \`etale lifting across the square-zero extension
\[
    \pi_0(A')
    \longrightarrow
    \pi_0(A)
\]
is unique up to unique isomorphism, and spectral \`etale algebras over a
connective base are determined by their classical truncations up to a
contractible space of choices. It follows that the \(\infty\)-groupoid
of pairs
\[
    (B',\alpha)
\]
with
\[
    \alpha
    \colon
    B'\otimes_{A'}A
    \xrightarrow{\simeq}
    B
\]
is contractible.
\end{proof}

\begin{corollary}[Local lifting of smooth algebras]
\label{cor:local-lifting-smooth-algebras}
In the situation of \Ref{thm:affine-cotangent-obstruction}, if $A\to B$ is
smooth, then $B$ admits a deformation over $A'$ after restriction to a
Zariski cover of $\Spec(B)$.
\end{corollary}

\begin{proof}
Zariski-locally, $L_{B/A}$ is finite free and $A\to B$ admits the standard
factorization of \Ref{thm:smooth-local-standard-form}. The brave new
affine-space factor lifts canonically by base change, and the \`etale
factor lifts uniquely by \Ref{cor:lifting-etale-algebras}. Their composite
gives the required local lift.
\end{proof}

\subsection{Lifting across general closed immersions}

\begin{theorem}[Local lifting across closed immersions]
\label{thm:infinitesimal-lifting-smooth-etale}
Let
\[
    i
    \colon
    Z
    \hookrightarrow
    S
\]
be a closed immersion of spectral schemes, and let
\[
    X
    \longrightarrow
    Z
\]
be smooth, respectively \`etale. Then every point of $X$ admits an open
neighbourhood $U\subseteq X$ for which there exists a smooth,
respectively \`etale, spectral scheme $\widetilde U\to S$ and an
equivalence
\[
    \widetilde U\times_SZ
    \simeq
    U.
\]
For an infinitesimal or nilpotent thickening the \`etale lift is unique
up to a contractible space of choices. No such uniqueness is asserted for
an arbitrary closed immersion.
\end{theorem}

\begin{proof}
The assertion is local on $S$, $Z$, and $X$, so fix a point and shrink
until all three schemes are affine.

Suppose first that $X\to Z$ is \`etale. Its classical truncation is an
ordinary \`etale morphism. The ordinary local lifting theorem for
\`etale morphisms across the closed immersion
\[
    Z_{\mathrm{cl}}
    \hookrightarrow
    S_{\mathrm{cl}}
\]
produces, after shrinking around the chosen point, an ordinary \`etale
scheme
\[
    \widetilde U_0
    \longrightarrow
    S_{\mathrm{cl}}
\]
whose base change to $Z_{\mathrm{cl}}$ is $U_{\mathrm{cl}}$. By
\Ref{thm:etale-invariance-classical-truncation}, $\widetilde U_0$ has a
canonical spectral \`etale lift
\[
    \widetilde U
    \longrightarrow
    S.
\]
The two \`etale $Z$-schemes $\widetilde U\times_SZ$ and $U$ have
equivalent classical truncations. Applying the same invariance theorem
over $Z$ gives
\[
    \widetilde U\times_SZ
    \simeq
    U.
\]
This proves local existence. Different choices of the ordinary local lift
need not be equivalent, which is why no uniqueness statement is made for
a general closed immersion. If $Z\hookrightarrow S$ is an infinitesimal or nilpotent
thickening, restriction induces an equivalence between the
$\infty$-categories of \`etale spectral schemes over $S$ and over $Z$.
Indeed, classical \`etale schemes are invariant under nilpotent
thickenings, and
\Ref{thm:etale-invariance-classical-truncation}
identifies the spectral \`etale categories with their classical
counterparts. Consequently the space of \`etale lifts is contractible.

Now suppose that $X\to Z$ is smooth. By
\Ref{thm:smooth-local-standard-form}, after shrinking around the chosen
point there is a factorization
\[
    U
    \longrightarrow
    \mathbb A_Z^n
    \longrightarrow
    Z
\]
whose first morphism is \`etale. The affine-space factor lifts
canonically to $\mathbb A_S^n\to S$. Apply the \`etale case to the closed
immersion
\[
    \mathbb A_Z^n
    \simeq
    \mathbb A_S^n\times_SZ
    \hookrightarrow
    \mathbb A_S^n.
\]
The resulting \`etale lift over $\mathbb A_S^n$, composed with
$\mathbb A_S^n\to S$, is smooth and restricts to $U\to Z$. See
\cite{LurieSAG,KhanBraveNewMotivicI}.
\end{proof}

\subsection{Regularity of closed immersions}

\begin{theorem}[Closed immersion between smooth schemes]
\label{thm:closed-between-smooth-regular}
Let $S$ be a spectral scheme. If
\[
    i
    \colon
    Z
    \hookrightarrow
    X
\]
is a closed immersion and both $Z\to S$ and $X\to S$ are smooth, then
$i$ is a regular closed immersion.
\end{theorem}

\begin{proof}
The assertion is local on $S$, $Z$, and $X$, so work with connective
commutative ring spectra
\[
    A
    \longrightarrow
    B
    \longrightarrow
    C,
\]
where $A\to B$ and $A\to C$ are smooth and $B\to C$ is surjective on
$\pi_0$.

The $A$-algebras $B$ and $C$ are compact. Applying
\Ref{lem:compact-arrow-undercategory} to $B\to C$ in
$\CAlg(\Sp)_{A/}$ shows that $C$ is compact as a $B$-algebra. Hence the
closed immersion is locally of finite presentation.

The transitivity triangle is
\[
    L_{B/A}\otimes_BC
    \longrightarrow
    L_{C/A}
    \longrightarrow
    L_{C/B}.
\]
Put
\[
    P:=L_{B/A}\otimes_BC,
    \qquad
    Q:=L_{C/A}.
\]
Both $P$ and $Q$ are finite locally free $C$-modules. Since
$\pi_0(B)\to\pi_0(C)$ is surjective,
\[
    \pi_0(L_{C/B})
    \simeq
    \Omega^1_{\pi_0(C)/\pi_0(B)}
    =
    0.
\]
Taking $\pi_0$ of the transitivity triangle gives a surjection
\[
    \pi_0(P)
    \longrightarrow
    \pi_0(Q).
\]
After a Zariski localization, choose a section of this surjection.

Because $Q$ is finite projective, the canonical map
\[
    \pi_0\Map_{\Mod_C}(Q,P)
    \longrightarrow
    \Hom_{\pi_0(C)}
    \bigl(
        \pi_0(Q),
        \pi_0(P)
    \bigr)
\]
is surjective. Hence the section lifts to a morphism
\[
    s
    \colon
    Q
    \longrightarrow
    P.
\]
The composite
\[
    Q
    \xrightarrow{s}
    P
    \longrightarrow
    Q
\]
is the identity on $\pi_0$. Both modules are flat, so the flat
homotopy-group formula shows that this composite is an equivalence.
Consequently
\[
    P
    \simeq
    K\oplus Q,
\]
where
\[
    K
    :=
    \operatorname{fib}(P\to Q)
\]
is finite locally free. Rotating the transitivity triangle gives
\[
    K
    \simeq
    L_{C/B}[-1]
    =
    N^*_{Z/X}.
\]
Thus the conormal module is finite locally free, and the closed immersion
is regular.
\end{proof}

\subsection{\'{E}tale neighbourhood of a smooth section}

\begin{theorem}[Spectral tubular neighbourhood]
\label{thm:spectral-tubular-neighbourhood}
Let
\[
    p
    \colon
    X
    \longrightarrow
    S
\]
be a smooth affine morphism, and let
\[
    s
    \colon
    S
    \longrightarrow
    X
\]
be a section. Then $s$ is a regular closed immersion. Put
\[
    N^*_{S/X}
    :=
    L_{S/X}[-1].
\]
There exists a Zariski open cover
\[
    \{S_\alpha\to S\}_\alpha
\]
such that, writing
\[
    X_\alpha
    :=
    X\times_SS_\alpha,
\]
there is an open neighbourhood
\[
    X_\alpha^\circ
    \subseteq
    X_\alpha
\]
of $s(S_\alpha)$ and an \`etale morphism over $S_\alpha$
\[
    \varphi_\alpha
    \colon
    X_\alpha^\circ
    \longrightarrow
    \mathbb V_{S_\alpha}
    \bigl(
        N^*_{S/X}|_{S_\alpha}
    \bigr)
\]
which carries the section to the zero section.

Under the canonical identifications
\[
    N^*_{S_\alpha/
    \mathbb V_{S_\alpha}
    (N^*_{S/X}|_{S_\alpha})}
    \simeq
    N^*_{S/X}|_{S_\alpha}
\]
and
\[
    N^*_{S_\alpha/X_\alpha^\circ}
    \simeq
    N^*_{S/X}|_{S_\alpha},
\]
the morphism on conormal modules induced by \(\varphi_\alpha\)
is homotopic to the identity of
\[
    N^*_{S/X}|_{S_\alpha}.
\]

If $S$ is affine, the cover may be taken to consist of $S$ itself.
\end{theorem}

\begin{proof}
We first verify directly that
\[
    s
    \colon
    S
    \longrightarrow
    X
\]
is a closed immersion. The assertion is local on \(S\). Write
\[
    S=\Spec(A),
    \qquad
    X=\Spec(B).
\]
The section is represented by an augmentation
\[
    \epsilon
    \colon
    B
    \longrightarrow
    A
\]
whose composite with the structural morphism
\[
    A
    \longrightarrow
    B
\]
is the identity. Consequently,
\[
    \pi_0(B)
    \longrightarrow
    \pi_0(A)
\]
is surjective. By the affine closed-immersion criterion,
\[
    s
    \colon
    S
    \hookrightarrow
    X
\]
is a closed immersion. It is regular by
\Ref{thm:closed-between-smooth-regular}, applied to the smooth
morphisms
\[
    X\longrightarrow S
    \qquad\text{and}\qquad
    S\longrightarrow S.
\]

The remaining assertion is Zariski-local on \(S\). Choose an affine open
cover
\[
    \{S_\alpha\to S\}_\alpha.
\]
After base change to one member of this cover and suppressing the index,
we may suppose that \(S\) is affine. Since \(p\) is affine, \(X\) is
then affine. Put
\[
    N
    :=
    N^*_{S/X}.
\]

Write
\[
    X
    \simeq
    \Spec_S(\mathcal A)
\]
with augmentation
\[
    \epsilon
    \colon
    \mathcal A
    \longrightarrow
    \mathcal O_S
\]
corresponding to \(s\). On classical truncations, put
\[
    \mathcal I
    :=
    \ker
    \left(
        \pi_0(\mathcal A)
        \longrightarrow
        \mathcal O_{S_{\mathrm{cl}}}
    \right).
\]

Transitivity for
\[
    S
    \xrightarrow{s}
    X
    \xrightarrow{p}
    S
\]
gives a canonical equivalence
\[
    N
    =
    L_{S/X}[-1]
    \simeq
    s^*L_{X/S}.
\]
The module \(L_{X/S}\) is finite locally free. The flat
homotopy-group formula, together with the degree-zero comparison for the
cotangent complex, therefore gives
\[
    \pi_0(N)
    \simeq
    \mathcal O_{S_{\mathrm{cl}}}
    \otimes_{\pi_0(\mathcal A)}
    \Omega^1_{\pi_0(\mathcal A)/
    \mathcal O_{S_{\mathrm{cl}}}}.
\]

By
\Ref{cor:classical-truncation-smooth}, the classical morphism
\[
    p_{\mathrm{cl}}
    \colon
    X_{\mathrm{cl}}
    \longrightarrow
    S_{\mathrm{cl}}
\]
is smooth. Its section
\[
    s_{\mathrm{cl}}
    \colon
    S_{\mathrm{cl}}
    \longrightarrow
    X_{\mathrm{cl}}
\]
is therefore a classical regular closed immersion. The classical
conormal identification for a section of a smooth morphism gives a
canonical isomorphism
\[
    \mathcal I/\mathcal I^2
    \xrightarrow{\simeq}
    \mathcal O_{S_{\mathrm{cl}}}
    \otimes_{\pi_0(\mathcal A)}
    \Omega^1_{\pi_0(\mathcal A)/
    \mathcal O_{S_{\mathrm{cl}}}}.
\]
Consequently,
\[
    \pi_0(N)
    \simeq
    \mathcal I/\mathcal I^2.
\]

The natural epimorphism
\[
    \mathcal I
    \longrightarrow
    \mathcal I/\mathcal I^2
\]
admits an
\(\mathcal O_{S_{\mathrm{cl}}}\)-linear section because \(S\) is affine
and \(\mathcal I/\mathcal I^2\) is finite projective. Composing such a
section with the preceding equivalence gives a morphism
\[
    \pi_0(N)
    \longrightarrow
    \mathcal I
    \hookrightarrow
    \pi_0(p_*\mathcal O_X).
\]

Since \(N\) is finite projective, the canonical map
\[
    \pi_0
    \Map_{\QCoh(S)}
    \bigl(
        N,
        p_*\mathcal O_X
    \bigr)
    \longrightarrow
    \Hom_{\mathcal O_{S_{\mathrm{cl}}}}
    \left(
        \pi_0(N),
        \pi_0(p_*\mathcal O_X)
    \right)
\]
is surjective. Choose a lift
\[
    \alpha
    \colon
    N
    \longrightarrow
    p_*\mathcal O_X.
\]

The composite
\[
    N
    \xrightarrow{\alpha}
    p_*\mathcal O_X
    \xrightarrow{\epsilon}
    \mathcal O_S
\]
induces the zero morphism on \(\pi_0\). Since maps out of a finite
projective module are detected on \(\pi_0\), this composite represents
the zero connected component of
\[
    \Map_{\QCoh(S)}(N,\mathcal O_S).
\]
Choose a nullhomotopy of this composite.

By the free-algebra adjunction, \(\alpha\), together with the chosen
nullhomotopy, determines a morphism of augmented quasi-coherent
commutative algebras
\[
    \operatorname{Sym}_{\mathcal O_S}(N)
    \longrightarrow
    p_*\mathcal O_X,
\]
and hence an \(S\)-morphism
\[
    q
    \colon
    X
    \longrightarrow
    \mathbb V_S(N)
\]
which carries \(s\) to the zero section.

The resulting morphism of pairs
\[
    (X,S)
    \longrightarrow
    \bigl(
        \mathbb V_S(N),S
    \bigr)
\]
induces a morphism of conormal modules
\[
    N^*_{S/\mathbb V_S(N)}
    \longrightarrow
    N^*_{S/X}.
\]
Under the canonical identifications
\[
    N^*_{S/\mathbb V_S(N)}
    \simeq
    N
    \simeq
    N^*_{S/X},
\]
the induced map on \(\pi_0\) is the identity. Indeed, it is the composite
obtained from the chosen section
\[
    \mathcal I/\mathcal I^2
    \longrightarrow
    \mathcal I
\]
followed by the quotient
\[
    \mathcal I
    \longrightarrow
    \mathcal I/\mathcal I^2.
\]

Both conormal modules are finite projective. A morphism between finite
projective modules which is an isomorphism on \(\pi_0\) is an
equivalence. Hence the induced conormal morphism is an equivalence and,
under the displayed identifications, is homotopic to the identity of
\(N\).

Let
\[
    Z
    :=
    X
    \times_{\mathbb V_S(N)}
    S
\]
be the zero fibre. Since
\[
    q\circ s
\]
is the zero section, \(s\) factors canonically as
\[
    S
    \longrightarrow
    Z
    \longrightarrow
    X.
\]

The closed immersion
\[
    Z
    \longrightarrow
    X
\]
is the base change of the zero section
\[
    S
    \longrightarrow
    \mathbb V_S(N).
\]
Cotangent base change therefore identifies its conormal module with the
pullback of \(N\). After restriction along
\[
    S
    \longrightarrow
    Z,
\]
the conormal morphism for the composable closed immersions
\[
    S
    \longrightarrow
    Z
    \longrightarrow
    X
\]
is precisely the conormal equivalence constructed above. The
transitivity triangle
\[
    L_{Z/X}|_S
    \longrightarrow
    L_{S/X}
    \longrightarrow
    L_{S/Z}
\]
therefore has its first morphism an equivalence. Hence
\[
    L_{S/Z}
    \simeq
    0.
\]

The morphism
\[
    S
    \longrightarrow
    Z
\]
is locally of finite presentation. Indeed, on an affine chart write
\[
    S=\Spec(A),
    \qquad
    X=\Spec(B),
    \qquad
    \mathbb V_S(N)=\Spec(R).
\]
Then
\[
    Z
    =
    \Spec(C),
    \qquad
    C
    :=
    B\otimes_RA.
\]
The \(A\)-algebras \(B\), \(R\), and \(A\) are compact. Compact objects
are closed under finite colimits, so \(C\) is compact in
\[
    \CAlg(\Sp)_{A/}.
\]
The morphism \(S\to Z\) is represented by the morphism
\[
    C
    \longrightarrow
    A.
\]
Applying
\Ref{lem:compact-arrow-undercategory}
to this morphism in
\(\CAlg(\Sp)_{A/}\) shows that \(A\) is compact as a \(C\)-algebra.
Thus \(S\to Z\) is locally of finite presentation. Together with
\[
    L_{S/Z}
    \simeq
    0,
\]
this shows that \(S\to Z\) is \`etale.

We now verify that
\[
    S
    \longrightarrow
    Z
\]
is a closed immersion. The morphism \(Z\to S\) is affine because \(Z\)
is a closed spectral subscheme of the affine \(S\)-scheme \(X\). The
morphism \(S\to Z\) is a section of \(Z\to S\). In the affine notation
above, it is represented by an augmented \(A\)-algebra
\[
    C
    \longrightarrow
    A
\]
whose composite
\[
    A
    \longrightarrow
    C
    \longrightarrow
    A
\]
is the identity. Consequently,
\[
    \pi_0(C)
    \longrightarrow
    \pi_0(A)
\]
is surjective. By the affine closed-immersion criterion,
\[
    S
    \longrightarrow
    Z
\]
is a closed immersion.

An \`etale closed immersion is an open and closed immersion. Therefore
there is a decomposition
\[
    Z
    \simeq
    S
    \coprod
    Z',
\]
where \(Z'\) is open and closed in \(Z\). Since
\[
    Z
    \hookrightarrow
    X
\]
is a closed immersion, the composite
\[
    Z'
    \hookrightarrow
    X
\]
is a closed immersion. Moreover, \(Z'\) is disjoint from \(s(S)\).

Replace \(X\) by the open complement
\[
    X_1
    :=
    X\setminus Z'.
\]
Then \(X_1\) is an open neighbourhood of \(s(S)\), and
\[
    X_1
    \times_{\mathbb V_S(N)}
    S
    \simeq
    S.
\]

Both \(X_1\) and \(\mathbb V_S(N)\) are locally of finite presentation
over \(S\). On affine charts over \(\Spec(A)\), their coordinate
algebras are compact \(A\)-algebras. The morphism
\[
    q|_{X_1}
    \colon
    X_1
    \longrightarrow
    \mathbb V_S(N)
\]
is therefore represented locally by a morphism between compact objects
of
\[
    \CAlg(\Sp)_{A/}.
\]
By
\Ref{lem:compact-arrow-undercategory},
it is locally of finite presentation.

Its relative cotangent complex fits into the transitivity triangle
\[
    q^*L_{\mathbb V_S(N)/S}
    \longrightarrow
    L_{X_1/S}
    \longrightarrow
    L_{X_1/\mathbb V_S(N)}.
\]
After pulling back along \(s\), the first two terms identify canonically
with \(N\), and the first morphism is the conormal equivalence already
identified, under these canonical identifications, with an endomorphism
homotopic to the identity. Hence
\[
    s^*L_{X_1/\mathbb V_S(N)}
    \simeq
    0.
\]

The relative cotangent complex
\[
    L_{X_1/\mathbb V_S(N)}
\]
is perfect, since it is the cofibre of a morphism between finite locally
free modules. The vanishing locus of a perfect quasi-coherent module is
open: on an affine chart, compactness implies that vanishing after
localization at a point already holds on a principal neighbourhood.
Therefore there is an open neighbourhood
\[
    X_0
    \subseteq
    X_1
\]
of \(s(S)\) such that
\[
    L_{X_0/\mathbb V_S(N)}
    \simeq
    0.
\]

The morphism
\[
    X_0
    \longrightarrow
    \mathbb V_S(N)
\]
is locally of finite presentation and has vanishing relative cotangent
complex. It is therefore \`etale. It carries \(s\) to the zero section.

Under the canonical identifications
\[
    N^*_{S/\mathbb V_S(N)}
    \simeq
    N
    \simeq
    N^*_{S/X_0},
\]
the induced morphism on conormal modules is homotopic to
\[
    \id_N.
\]

Restoring the affine-cover index gives, for every \(\alpha\), an open
neighbourhood
\[
    X_\alpha^\circ
    \subseteq
    X_\alpha
\]
and an \`etale morphism
\[
    \varphi_\alpha
    \colon
    X_\alpha^\circ
    \longrightarrow
    \mathbb V_{S_\alpha}
    \bigl(
        N^*_{S/X}|_{S_\alpha}
    \bigr)
\]
with all the asserted properties. When \(S\) is affine, the preceding
construction applies globally, so the cover may be taken to consist of
\(S\) itself.
\end{proof}

\section{The smooth spectral site over a base}
\label{sec:smooth-spectral-site}

Fix for the remainder of the chapter a quasi-compact quasi-separated
spectral scheme
\[
    S.
\]

\subsection{Smooth spectral schemes of finite presentation}

\begin{definition}[The smooth spectral test category]
\label{def:smooth-spectral-test-category}
Let
\[
    \operatorname{Sm}^{\mathrm{fp}}_{/S}
\]
denote the full subcategory of spectral schemes over $S$ spanned by the
morphisms
\[
    X
    \longrightarrow
    S
\]
which are smooth and of finite presentation in the sense of
\Ref{def:finitely-presented-spectral-morphism}.
\end{definition}

Smoothness is an affine-local condition and therefore supplies local
finite presentation. The superscript $\mathrm{fp}$ additionally imposes
the global quasi-compactness and quasi-separatedness conditions required
for the motivic test category.

\begin{proposition}[Essential smallness]
\label{prop:smooth-site-essentially-small}
The $\infty$-category
\[
    \operatorname{Sm}^{\mathrm{fp}}_{/S}
\]
is essentially small in a universe containing $S$.
\end{proposition}

\begin{proof}
Choose a finite affine open cover
\[
    S
    =
    \bigcup_{a=1}^r
    S_a,
    \qquad
    S_a=\Spec(A_a),
\]
which exists because $S$ is quasi-compact.

Let $X\to S$ be smooth and of finite presentation. Since the morphism is
quasi-compact, every pullback
\[
    X_a
    :=
    X\times_SS_a
\]
is quasi-compact. Since the morphism is quasi-separated and $S_a$ is
quasi-separated, each $X_a$ is quasi-separated. Consequently each $X_a$
admits a finite affine open cover
\[
    X_a
    =
    \bigcup_{b=1}^{m_a}
    \Spec(B_{ab}).
\]
After refining these finite affine covers, local finite presentation
ensures that every $B_{ab}$ is compact in
\[
    \CAlg(\Sp)_{A_a/}.
\]

For a fixed connective commutative ring spectrum $A_a$, the compact
objects of
\[
    \CAlg(\Sp)_{A_a/}
\]
form an essentially small $\infty$-category. Hence there is only an
essentially small collection of possible affine charts
$\Spec(B_{ab})$.

Because $X$ is quasi-separated, intersections of the finitely many affine
charts are quasi-compact. They therefore admit finite affine covers by
compact algebras over the corresponding base charts. The gluing maps,
their iterated compatibility maps, and all higher coherence data belong
to mapping spaces between objects of essentially small
$\infty$-categories. There is therefore only a set of equivalence classes
of such finite affine descent presentations in the chosen universe.

Every object of $\operatorname{Sm}^{\mathrm{fp}}_{/S}$ admits such a
presentation. Hence the category is essentially small.
\end{proof}

\begin{convention}[Small smooth site]
\label{conv:small-smooth-site}
The essentially $\mathbb U$-small category
\[
    \operatorname{Sm}^{\mathrm{fp}}_{/S}
\]
is used without choosing a skeleton. Its presheaf and sheaf categories are
formed in the larger universe $\mathbb V$, as in
\Ref{conv:size-spectral-ag}. In particular, geometric base change remains
a canonical functor on the stated categories.
\end{convention}

\begin{proposition}[Finite products]
\label{prop:smooth-site-finite-products}
The category $\operatorname{Sm}^{\mathrm{fp}}_{/S}$ admits finite products.
The terminal object is $S$, and the product of $X\to S$ and $Y\to S$ is
the spectral fibre product
\[
    X\times_SY.
\]
\end{proposition}\begin{proof}
The terminal morphism
\[
    S\longrightarrow S
\]
is smooth and of finite presentation.

Fibre products of spectral schemes exist by
\Ref{thm:spectral-schemes-fibre-products}. Let $X\to S$ and $Y\to S$ be
objects of $\operatorname{Sm}^{\mathrm{fp}}_{/S}$. The projection
\[
    X\times_SY
    \longrightarrow
    Y
\]
is obtained from $X\to S$ by base change and is therefore smooth and of
finite presentation by
\Ref{prop:global-smooth-etale-permanence} and
\Ref{prop:global-finite-presentation-permanence}. Composing with
$Y\to S$ shows that
\[
    X\times_SY
    \longrightarrow
    S
\]
is smooth and of finite presentation.
\end{proof}

\begin{proposition}[Finite coproducts]
\label{prop:smooth-site-finite-coproducts}
The category $\operatorname{Sm}^{\mathrm{fp}}_{/S}$ admits finite
coproducts. The initial object is the empty spectral scheme.
\end{proposition}

\begin{proof}
Finite coproducts of spectral schemes exist by
\Ref{prop:spectral-schemes-coproducts}. Smoothness and finite presentation
are checked separately on components.
\end{proof}

\subsection{Functoriality in the base}

\begin{proposition}[Geometric base change]
\label{prop:smooth-site-geometric-base-change}
Every morphism of spectral schemes
\[
    f
    \colon
    T
    \longrightarrow
    S
\]
induces a functor
\[
    f^*_{\mathrm{geom}}
    \colon
    \operatorname{Sm}^{\mathrm{fp}}_{/S}
    \longrightarrow
    \operatorname{Sm}^{\mathrm{fp}}_{/T},
    \qquad
    X
    \longmapsto
    X\times_ST.
\]
For composable morphisms
\[
    R
    \xrightarrow{g}
    T
    \xrightarrow{f}
    S,
\]
there is a canonical coherent equivalence
\[
    (f\circ g)^*_{\mathrm{geom}}
    \simeq
    g^*_{\mathrm{geom}}f^*_{\mathrm{geom}}.
\]
\end{proposition}

\begin{proof}
Let
\[
    X
    \longrightarrow
    S
\]
be smooth and of finite presentation. Smoothness is stable under
arbitrary base change by
\Ref{prop:global-smooth-etale-permanence}, and finite presentation is
stable under arbitrary base change by
\Ref{prop:global-finite-presentation-permanence}. Hence
\[
    X\times_ST
    \longrightarrow
    T
\]
belongs to
\[
    \operatorname{Sm}^{\mathrm{fp}}_{/T}.
\]

For composable morphisms
\[
    R
    \xrightarrow{g}
    T
    \xrightarrow{f}
    S,
\]
the canonical associativity equivalence for fibre products gives
\[
    (X\times_ST)\times_TR
    \simeq
    X\times_SR.
\]
These equivalences are natural in \(X\) and supply
\[
    (f\circ g)^*_{\mathrm{geom}}
    \simeq
    g^*_{\mathrm{geom}}f^*_{\mathrm{geom}}.
\]
Their higher coherence is inherited from the universal property of
limits.
\end{proof}

\begin{proposition}[Smooth composition functor]
\label{prop:smooth-composition-functor}
Let
\[
    p
    \colon
    X
    \longrightarrow
    S
\]
be an object of $\operatorname{Sm}^{\mathrm{fp}}_{/S}$. Composition with
$p$ defines a functor
\[
    p_{\sharp}^{\mathrm{geom}}
    \colon
    \operatorname{Sm}^{\mathrm{fp}}_{/X}
    \longrightarrow
    \operatorname{Sm}^{\mathrm{fp}}_{/S}.
\]
\end{proposition}\begin{proof}
Let $Y\to X$ be smooth and of finite presentation. Since
\[
    X\longrightarrow S
\]
is also smooth and of finite presentation, their composite
\[
    Y
    \longrightarrow
    X
    \longrightarrow
    S
\]
is smooth by
\Ref{prop:global-smooth-etale-permanence} and of finite presentation by
\Ref{prop:global-finite-presentation-permanence}. Thus composition with
$p$ defines the asserted functor.
\end{proof}

\subsection{The brave new interval object}

\begin{proposition}[The affine line belongs to the smooth site]
\label{prop:affine-line-in-smooth-site}
The projection
\[
    \mathbb A^1_S
    \longrightarrow
    S
\]
is an object of $\operatorname{Sm}^{\mathrm{fp}}_{/S}$.
\end{proposition}

\begin{proof}
This is the case $n=1$ of
\Ref{prop:brave-new-affine-space-smooth}.
\end{proof}

\begin{construction}[Interval structure maps]
\label{con:affine-line-interval-maps}
The zero and unit sections of
\Ref{con:zero-unit-sections-affine-line} give
\[
    0,1
    \colon
    S
    \rightrightarrows
    \mathbb A^1_S.
\]

The multiplication morphism
\[
    \mu
    \colon
    \mathbb A^1_S
    \times_S
    \mathbb A^1_S
    \longrightarrow
    \mathbb A^1_S
\]
is induced contravariantly by the morphism of quasi-coherent
commutative algebras
\[
    \operatorname{Sym}_{\mathcal O_S}(\mathcal O_S)
    \longrightarrow
    \operatorname{Sym}_{\mathcal O_S}
    \bigl(
        \mathcal O_S\oplus\mathcal O_S
    \bigr)
\]
which carries the universal generator \(t\) to the product
\[
    t_1t_2
\]
of the two universal generators.

The morphism \(\mu\) is coherently associative and commutative, the
section \(1\) is a unit, and the section \(0\) is absorbing:
\[
    \mu\circ(1,\id)
    \simeq
    \id
    \simeq
    \mu\circ(\id,1),
\]
and
\[
    \mu\circ(0,\id)
    \simeq
    0
    \simeq
    \mu\circ(\id,0).
\]
Thus \(\mathbb A^1_S\) is a commutative monoid object over \(S\)
equipped with a distinguished absorbing section. These maps are
compatible with arbitrary base change in \(S\).
\end{construction}

\begin{proof}
The sections were constructed operadically in
\Ref{con:zero-unit-sections-affine-line}.

Affine-locally on \(S=\Spec(A)\), write
\[
    A\{t\}
    :=
    \operatorname{Sym}_A(A)
\]
and
\[
    A\{t_1,t_2\}
    :=
    \operatorname{Sym}_A(A\oplus A).
\]
The morphism
\[
    \mu
    \colon
    \mathbb A^1_S\times_S\mathbb A^1_S
    \longrightarrow
    \mathbb A^1_S
\]
is represented contravariantly by the \(A\)-algebra morphism
\[
    A\{t\}
    \longrightarrow
    A\{t_1,t_2\},
    \qquad
    t
    \longmapsto
    t_1t_2.
\]

Associativity and commutativity follow from the coherently associative
and commutative multiplication in the commutative algebra
\[
    A\{t_1,t_2,t_3\}.
\]
The identities
\[
    1\cdot t=t=t\cdot1
\]
and
\[
    0\cdot t=0=t\cdot0
\]
give the unit and absorption equivalences after passing
contravariantly to affine spectra. All coherence data are inherited
from the commutative operadic algebra structure.

Finally, the free commutative algebra construction and the universal
generators commute with extension of scalars. The construction of
\(\mu\), together with the zero and unit sections, is therefore
compatible with arbitrary base change by
\Ref{prop:affine-space-base-change}.
\end{proof}

\section{Spectral Nisnevich geometry and the brave new motivic site}
\label{sec:spectral-nisnevich-motivic-input}

\subsection{Nisnevich distinguished squares}

\begin{definition}[Spectral Nisnevich square]
\label{def:spectral-nisnevich-square}
A Cartesian square in $\operatorname{Sm}^{\mathrm{fp}}_{/S}$
\[
\begin{tikzcd}[column sep=large, row sep=large]
    W
        \arrow[r]
        \arrow[d]
    &
    V
        \arrow[d, "p"]
    \\
    U
        \arrow[r, "j"']
    &
    X
\end{tikzcd}
\]
is a \textbf{spectral Nisnevich distinguished square} if:
\begin{enumerate}[label=(\roman*)]
    \item $j:U\to X$ is an open immersion;

    \item $p:V\to X$ is \`etale;

    \item there exists a closed immersion
    \[
        i
        \colon
        Z
        \hookrightarrow
        X
    \]
    together with an equivalence over $X$
    \[
        U
        \simeq
        X\setminus Z
    \]
    such that the base-change morphism
    \[
        V\times_XZ
        \longrightarrow
        Z
    \]
    is an equivalence.
\end{enumerate}
\end{definition}

The square may be represented together with a complementary closed
subscheme by
\[
\begin{tikzcd}[column sep=large, row sep=large]
    W=V\times_XU
        \arrow[r]
        \arrow[d]
    &
    V
        \arrow[d, "p"]
    &
    V\times_XZ
        \arrow[l]
        \arrow[d, "\simeq"]
    \\
    U
        \arrow[r, "j"']
    &
    X
    &
    Z
        \arrow[l, hook'].
\end{tikzcd}
\]

\begin{proposition}[Base change of Nisnevich squares]
\label{prop:nisnevich-square-base-change}
Spectral Nisnevich distinguished squares are stable under arbitrary base
change in $\operatorname{Sm}^{\mathrm{fp}}_{/S}$ and under arbitrary
change of the base spectral scheme $S$.
\end{proposition}

\begin{proof}
Let
\[
\begin{tikzcd}
    W \arrow[r] \arrow[d]
    &
    V \arrow[d, "p"]
    \\
    U \arrow[r, "j"']
    &
    X
\end{tikzcd}
\]
be distinguished, and choose a complementary closed immersion
\[
    Z\hookrightarrow X
\]
as in \Ref{def:spectral-nisnevich-square}. Let $X'\to X$ be any morphism
of spectral schemes and form the base-changed square.

Open immersions, \`etale morphisms, and closed immersions are stable under
base change. By \Ref{prop:complement-base-change}, there is a canonical
equivalence
\[
    X'\times_XU
    \simeq
    X'
    \setminus
    (X'\times_XZ).
\]
Moreover,
\[
\begin{aligned}
    (V\times_XX')
    \times_{X'}
    (Z\times_XX')
    &\simeq
    (V\times_XZ)\times_XX'
    \\
    &\simeq
    Z\times_XX'.
\end{aligned}
\]
Thus the base-changed square is distinguished. The same argument applies
to change of the ambient base spectral scheme.
\end{proof}

\begin{proposition}[Classical detection of the complementary condition]
\label{prop:nisnevich-classical-detection}
Consider a Cartesian square
\[
\begin{tikzcd}
    W \arrow[r] \arrow[d]
    &
    V \arrow[d, "p"]
    \\
    U \arrow[r, "j"']
    &
    X
\end{tikzcd}
\]
in $\operatorname{Sm}^{\mathrm{fp}}_{/S}$. Suppose that $j$ is an open
immersion and $p$ is \`etale. Then the square is a spectral Nisnevich
distinguished square if and only if its classical truncation is a
Nisnevich distinguished square of ordinary schemes.
\end{proposition}

\begin{proof}
Suppose first that the spectral square is distinguished, and choose a
complementary closed immersion
\[
    Z\hookrightarrow X.
\]
Classical truncation carries its open complement to the ordinary open
complement:
\[
    U_{\mathrm{cl}}
    \simeq
    X_{\mathrm{cl}}\setminus Z_{\mathrm{cl}}.
\]
The morphism $p$ is \`etale and therefore flat. Hence classical
truncation commutes with the base change along $p$, giving
\[
    (V\times_XZ)_{\mathrm{cl}}
    \simeq
    V_{\mathrm{cl}}
    \times_{X_{\mathrm{cl}}}
    Z_{\mathrm{cl}}.
\]
The equivalence $V\times_XZ\simeq Z$ therefore yields the ordinary
Nisnevich complementary equivalence.

Conversely, suppose that the classical square is Nisnevich
distinguished. Let
\[
    Z_{\mathrm{red}}
    :=
    \bigl(
        X_{\mathrm{cl}}
        \setminus
        U_{\mathrm{cl}}
    \bigr)_{\mathrm{red}}
\]
and regard $Z_{\mathrm{red}}$ as a classical spectral scheme equipped
with the composite closed immersion
\[
    Z_{\mathrm{red}}
    \hookrightarrow
    X_{\mathrm{cl}}
    \longrightarrow
    X.
\]
By the equivalence between open spectral subschemes of $X$ and open
subschemes of $X_{\mathrm{cl}}$, there is an equivalence
\[
    U
    \simeq
    X\setminus Z_{\mathrm{red}}.
\]

The base-change morphism
\[
    q
    \colon
    V\times_XZ_{\mathrm{red}}
    \longrightarrow
    Z_{\mathrm{red}}
\]
is \`etale. By the ordinary Nisnevich condition, its classical
truncation
\[
    q_{\mathrm{cl}}
    \colon
    (V\times_XZ_{\mathrm{red}})_{\mathrm{cl}}
    \longrightarrow
    Z_{\mathrm{red}}
\]
is an equivalence.

By
\Ref{thm:etale-invariance-classical-truncation},
an \`etale spectral scheme over $Z_{\mathrm{red}}$ is determined, up to
a contractible space of choices, by its classical truncation. Therefore
$q$ is an equivalence. Thus the spectral square is distinguished.
\end{proof}

\begin{proposition}[Pointwise Nisnevich criterion]
\label{prop:nisnevich-pointwise-criterion}
Consider a Cartesian square in $\operatorname{Sm}^{\mathrm{fp}}_{/S}$
\[
\begin{tikzcd}
    W \arrow[r] \arrow[d]
    &
    V \arrow[d,"p"]
    \\
    U \arrow[r,"j"']
    &
    X
\end{tikzcd}
\]
with $j$ an open immersion and $p$ \`etale. Then the square is
Nisnevich distinguished if and only if, for every point
\[
    x
    \in
    |X|\setminus|U|,
\]
there is a unique point $v\in|V|$ above $x$ and the induced map of
residue fields
\[
    \kappa(x)
    \longrightarrow
    \kappa(v)
\]
is an isomorphism.
\end{proposition}

\begin{proof}
By \Ref{prop:nisnevich-classical-detection}, the assertion is equivalent
to the corresponding statement for the classical truncation of the
square. The ordinary pointwise criterion for Nisnevich distinguished
squares identifies the stated residue-field condition with the assertion
that the \`etale morphism is an isomorphism over the reduced closed
complement. Translating back through classical truncation proves the
spectral statement.
\end{proof}

\subsection{The spectral Nisnevich topology}

\begin{definition}[Spectral Nisnevich topology]
\label{def:spectral-nisnevich-topology}
The \textbf{spectral Nisnevich topology}
\[
    \tau_{\mathrm{Nis}}
\]
on $\operatorname{Sm}^{\mathrm{fp}}_{/S}$ is the Grothendieck topology
generated by:
\begin{enumerate}[label=(\roman*)]
    \item the empty covering family of the initial object;

    \item the two-member covering families
    \[
        \{U\to X,\,V\to X\}
    \]
    arising from spectral Nisnevich distinguished squares.
\end{enumerate}
\end{definition}

\begin{proposition}[Pullback stability of Nisnevich generators]
\label{prop:nisnevich-pretopology}
The two-member families associated to spectral Nisnevich distinguished
squares are stable under pullback. Together with the empty covering
family of the initial object, they therefore generate the Grothendieck
topology of \Ref{def:spectral-nisnevich-topology}.
\end{proposition}

\begin{proof}
Pullback stability is \Ref{prop:nisnevich-square-base-change}. Any
pullback-stable collection of covering families determines a smallest
Grothendieck topology containing those families. By definition, this
smallest topology is $\tau_{\mathrm{Nis}}$.
\end{proof}

\begin{definition}[Spectral \`etale topology on the smooth site]
\label{def:spectral-etale-topology-smooth-site}
The \textbf{spectral \`etale topology}
\[
    \tau_{\mathrm{\acute et}}
\]
on
\[
    \operatorname{Sm}^{\mathrm{fp}}_{/S}
\]
is the Grothendieck topology generated by families
\[
    \{X_\alpha\to X\}_{\alpha\in I}
\]
such that:
\begin{enumerate}[label=(\roman*)]
    \item every morphism
    \[
        X_\alpha
        \longrightarrow
        X
    \]
    is \`etale;

    \item the induced map of underlying topological spaces
    \[
        \coprod_{\alpha\in I}|X_\alpha|
        \longrightarrow
        |X|
    \]
    is surjective.
\end{enumerate}
Equivalently, this is the topology induced on
\(\operatorname{Sm}^{\mathrm{fp}}_{/S}\) from the big spectral
\`etale site over \(S\).
\end{definition}

\begin{theorem}[Subcanonicity of the spectral Nisnevich topology]
\label{thm:spectral-nisnevich-subcanonical}
Every representable presheaf on
\(\operatorname{Sm}^{\mathrm{fp}}_{/S}\) is a sheaf for
\(\tau_{\mathrm{Nis}}\).
\end{theorem}

\begin{proof}
Let
\[
    \{U\to X,\;V\to X\}
\]
be a generating two-member family arising from a spectral Nisnevich
distinguished square
\[
\begin{tikzcd}[column sep=large, row sep=large]
    W
        \arrow[r]
        \arrow[d]
    &
    V
        \arrow[d, "p"]
    \\
    U
        \arrow[r, "j"']
    &
    X.
\end{tikzcd}
\]
The morphism
\[
    j
    \colon
    U
    \longrightarrow
    X
\]
is an open immersion and hence is \`etale by
\Ref{prop:open-immersions-etale}. The morphism
\[
    p
    \colon
    V
    \longrightarrow
    X
\]
is \`etale by the definition of a spectral Nisnevich distinguished
square.

Choose a complementary closed immersion
\[
    Z
    \hookrightarrow
    X
\]
such that
\[
    U
    \simeq
    X\setminus Z
\]
and
\[
    V\times_XZ
    \xrightarrow{\simeq}
    Z.
\]
The image of \(U\to X\) contains the open complement of \(Z\). The
equivalence
\[
    V\times_XZ
    \simeq
    Z
\]
shows that the image of \(V\to X\) contains every point of \(Z\).
Consequently,
\[
    |U|\cup|V|
    =
    |X|.
\]
Thus
\[
    \{U\to X,\;V\to X\}
\]
is a jointly surjective family of \`etale morphisms and hence is a
covering family for
\[
    \tau_{\mathrm{\acute et}}.
\]

Every generating covering family for
\(\tau_{\mathrm{Nis}}\) is therefore a
\(\tau_{\mathrm{\acute et}}\)-cover. Equivalently,
\[
    \tau_{\mathrm{Nis}}
    \subseteq
    \tau_{\mathrm{\acute et}},
\]
so the spectral Nisnevich topology is coarser than the spectral
\`etale topology.

Spectral schemes and their morphisms satisfy effective \`etale descent.
Consequently, for every
\[
    Y
    \in
    \operatorname{Sm}^{\mathrm{fp}}_{/S},
\]
the representable presheaf
\[
    h_S(Y)
    \colon
    (\operatorname{Sm}^{\mathrm{fp}}_{/S})^{\op}
    \longrightarrow
    \mathcal S,
    \qquad
    h_S(Y)(T)
    =
    \Map_S(T,Y),
\]
is a sheaf for
\(\tau_{\mathrm{\acute et}}\).

A sheaf for a finer topology is a sheaf for every coarser topology.
Therefore every representable presheaf is a sheaf for
\[
    \tau_{\mathrm{Nis}}.
\]
Hence the spectral Nisnevich topology is subcanonical. See
\cite{LurieSAG,KhanBraveNewMotivicI}.
\end{proof}

\subsection{Excision}

\begin{definition}[Nisnevich-excisive presheaf]
\label{def:nisnevich-excisive-presheaf}
A presheaf
\[
    F
    \colon
    (\operatorname{Sm}^{\mathrm{fp}}_{/S})^{\op}
    \longrightarrow
    \mathcal S
\]
is \textbf{Nisnevich-excisive} if:
\begin{enumerate}[label=(\roman*)]
    \item the space $F(\varnothing)$ is contractible;

    \item every spectral Nisnevich distinguished square is carried to a
    Cartesian square
    \[
    \begin{tikzcd}[column sep=large, row sep=large]
        F(X)
            \arrow[r]
            \arrow[d]
        &
        F(U)
            \arrow[d]
        \\
        F(V)
            \arrow[r]
        &
        F(W).
    \end{tikzcd}
    \]
\end{enumerate}
\end{definition}

\begin{theorem}[Excision criterion for Nisnevich descent]
\label{thm:nisnevich-excision-criterion}
A presheaf of spaces on $\operatorname{Sm}^{\mathrm{fp}}_{/S}$ is a sheaf
for $\tau_{\mathrm{Nis}}$ if and only if it is Nisnevich-excisive.
\end{theorem}

\begin{proof}
Let
\[
    E_{\mathrm{Nis}}
\]
denote the collection of spectral Nisnevich distinguished squares in
\[
    \operatorname{Sm}^{\mathrm{fp}}_{/S}.
\]

The empty spectral scheme is a strict initial object: every morphism
\[
    X
    \longrightarrow
    \varnothing
\]
is an equivalence. Moreover,
\Ref{prop:nisnevich-square-base-change}
shows that \(E_{\mathrm{Nis}}\) is stable under pullback. Hence
\(E_{\mathrm{Nis}}\) is an excision structure.

Consider a distinguished square
\[
\begin{tikzcd}
    W \arrow[r] \arrow[d]
    &
    V \arrow[d, "p"]
    \\
    U \arrow[r, "j"']
    &
    X.
\end{tikzcd}
\]

The morphism
\[
    j
    \colon
    U
    \longrightarrow
    X
\]
is an open immersion and hence a monomorphism. This verifies the
monomorphism condition for a topological excision structure.

It remains to verify stability under the diagonal construction. Consider
the induced Cartesian square
\[
\begin{tikzcd}[column sep=large, row sep=large]
    W
        \arrow[r]
        \arrow[d, "\Delta_{W/U}"']
    &
    V
        \arrow[d, "\Delta_{V/X}"]
    \\
    W\times_UW
        \arrow[r]
    &
    V\times_XV.
\end{tikzcd}
\]

Since
\[
    W
    \simeq
    U\times_XV,
\]
there is a canonical equivalence
\[
    W\times_UW
    \simeq
    U\times_X(V\times_XV).
\]
Consequently, the lower horizontal morphism
\[
    W\times_UW
    \longrightarrow
    V\times_XV
\]
is the base change of the open immersion
\[
    U
    \longrightarrow
    X
\]
and is therefore an open immersion.

By
\Ref{prop:etale-diagonal-and-monomorphisms},
the diagonal
\[
    \Delta_{V/X}
    \colon
    V
    \longrightarrow
    V\times_XV
\]
of the \`etale morphism \(p\) is an open immersion and, in particular,
is \`etale.

Choose a complementary closed immersion
\[
    Z
    \hookrightarrow
    X
\]
for the original distinguished square, so that
\[
    U
    \simeq
    X\setminus Z
\]
and
\[
    V\times_XZ
    \xrightarrow{\simeq}
    Z.
\]

The closed complement of
\[
    W\times_UW
    \longrightarrow
    V\times_XV
\]
is the base change
\[
    D
    :=
    (V\times_XV)\times_XZ
    \hookrightarrow
    V\times_XV.
\]
There are canonical equivalences
\[
\begin{aligned}
    D
    &\simeq
    (V\times_XZ)\times_Z(V\times_XZ)
    \\
    &\simeq
    Z\times_ZZ
    \\
    &\simeq
    Z.
\end{aligned}
\]

The restriction of the diagonal to this closed complement is the
morphism
\[
    V\times_{V\times_XV}D
    \longrightarrow
    D.
\]
Under the preceding identifications, this is the diagonal of
\[
    V\times_XZ
    \longrightarrow
    Z.
\]
Since
\[
    V\times_XZ
    \xrightarrow{\simeq}
    Z,
\]
that diagonal is an equivalence. Therefore the diagonal square is again
a spectral Nisnevich distinguished square.

Thus \(E_{\mathrm{Nis}}\) satisfies the monomorphism and diagonal
conditions and is a topological excision structure.

For a topological excision structure, the sheaves for the generated
Grothendieck topology are exactly the presheaves which carry the initial
object to a terminal space and every distinguished square to a Cartesian
square. Applying the topological excision theorem to
\(E_{\mathrm{Nis}}\) shows that a presheaf is a
\(\tau_{\mathrm{Nis}}\)-sheaf if and only if it is
Nisnevich-excisive. See
\cite{VoevodskyCDStructures,KhanBraveNewMotivicI}.
\end{proof}

\subsection{Representables and the affine-line projection class}

\begin{definition}[Smooth representable presheaf]
\label{def:smooth-representable-presheaf}
For $X\in\operatorname{Sm}^{\mathrm{fp}}_{/S}$, write
\[
    h_S(X)
    \colon
    (\operatorname{Sm}^{\mathrm{fp}}_{/S})^{\op}
    \longrightarrow
    \mathcal S
\]
for the representable presheaf
\[
    h_S(X)(T)
    :=
    \Map_S(T,X).
\]
\end{definition}

\begin{proposition}[Yoneda on the smooth site]
\label{prop:yoneda-smooth-site}
The assignment $X\mapsto h_S(X)$ defines a fully faithful functor
\[
    h_S
    \colon
    \operatorname{Sm}^{\mathrm{fp}}_{/S}
    \hookrightarrow
    \PSh
    \bigl(
        \operatorname{Sm}^{\mathrm{fp}}_{/S}
    \bigr).
\]
Every $h_S(X)$ is Nisnevich-excisive.
\end{proposition}

\begin{proof}
Full faithfulness is the Yoneda lemma. Nisnevich excision follows from
\Ref{thm:spectral-nisnevich-subcanonical} and
\Ref{thm:nisnevich-excision-criterion}.
\end{proof}

\begin{definition}[Brave new affine-line projection class]
\label{def:affine-line-projection-class}
Let
\[
    W_{\mathbb A^1,S}
\]
be the class of morphisms of presheaves
\[
    h_S(X\times_S\mathbb A^1_S)
    \longrightarrow
    h_S(X)
\]
induced by the projections
\[
    X\times_S\mathbb A^1_S
    \longrightarrow
    X,
\]
as $X$ ranges over $\operatorname{Sm}^{\mathrm{fp}}_{/S}$.
\end{definition}

\begin{proposition}[Pullback stability of affine-line projections]
\label{prop:affine-line-projection-pullback}
The class
\[
    W_{\mathbb A^1,S}
\]
is stable under pullback along morphisms between representable
presheaves.

More precisely, for every morphism
\[
    Y
    \longrightarrow
    X
\]
in
\[
    \operatorname{Sm}^{\mathrm{fp}}_{/S},
\]
the pullback of
\[
    h_S(X\times_S\mathbb A^1_S)
    \longrightarrow
    h_S(X)
\]
along
\[
    h_S(Y)
    \longrightarrow
    h_S(X)
\]
is canonically equivalent to
\[
    h_S(Y\times_S\mathbb A^1_S)
    \longrightarrow
    h_S(Y).
\]
In particular, the resulting pullback morphism again belongs to
\(W_{\mathbb A^1,S}\).
\end{proposition}

\begin{proof}
For every morphism
\[
    Y
    \longrightarrow
    X
\]
over \(S\), there is a canonical Cartesian square
\[
\begin{tikzcd}[column sep=large, row sep=large]
    Y\times_S\mathbb A^1_S
        \arrow[r]
        \arrow[d]
    &
    X\times_S\mathbb A^1_S
        \arrow[d]
    \\
    Y
        \arrow[r]
    &
    X.
\end{tikzcd}
\]
Indeed,
\[
\begin{aligned}
    Y\times_X
    \bigl(
        X\times_S\mathbb A^1_S
    \bigr)
    &\simeq
    Y\times_S\mathbb A^1_S.
\end{aligned}
\]

The Yoneda embedding preserves limits. Applying \(h_S\) to the displayed
Cartesian square therefore gives a Cartesian square of presheaves
\[
\begin{tikzcd}[column sep=large, row sep=large]
    h_S(Y\times_S\mathbb A^1_S)
        \arrow[r]
        \arrow[d]
    &
    h_S(X\times_S\mathbb A^1_S)
        \arrow[d]
    \\
    h_S(Y)
        \arrow[r]
    &
    h_S(X).
\end{tikzcd}
\]
Thus the pullback of the affine-line projection associated to \(X\) is
the affine-line projection associated to \(Y\), and hence again belongs
to
\[
    W_{\mathbb A^1,S}.
\]
\end{proof}

\begin{proposition}[Base change of the affine-line projection class]
\label{prop:affine-line-projection-base-change}
For every morphism of spectral schemes
\[
    f
    \colon
    T
    \longrightarrow
    S,
\]
put
\[
    f^{\mathrm{PSh}}_!
    :=
    \operatorname{Lan}_{(f^*_{\mathrm{geom}})^{\op}}
    \colon
    \PSh
    \bigl(
        \operatorname{Sm}^{\mathrm{fp}}_{/S}
    \bigr)
    \longrightarrow
    \PSh
    \bigl(
        \operatorname{Sm}^{\mathrm{fp}}_{/T}
    \bigr).
\]
There are canonical equivalences
\[
    f^{\mathrm{PSh}}_!h_S(X)
    \simeq
    h_T(X\times_ST).
\]
Under these equivalences, $f^{\mathrm{PSh}}_!$ carries the morphism
\[
    h_S(X\times_S\mathbb A^1_S)
    \longrightarrow
    h_S(X)
\]
to the morphism
\[
    h_T
    \bigl(
        (X\times_ST)\times_T\mathbb A^1_T
    \bigr)
    \longrightarrow
    h_T(X\times_ST).
\]
Consequently,
\[
    f^{\mathrm{PSh}}_!
    \bigl(
        W_{\mathbb A^1,S}
    \bigr)
    \subseteq
    W_{\mathbb A^1,T}.
\]
\end{proposition}

\begin{proof}
Left Kan extension along a functor carries representable presheaves to
the corresponding representable presheaves. Hence
\[
    f^{\mathrm{PSh}}_!h_S(X)
    \simeq
    h_T(X\times_ST).
\]
Use the canonical base-change equivalence
\[
    \mathbb A^1_S\times_ST
    \simeq
    \mathbb A^1_T
\]
from \Ref{prop:affine-space-base-change} and associativity of fibre
products.
\end{proof}

\subsection{The brave new motivic site theorem}

\begin{theorem}[Brave new motivic site package]
\label{thm:motivic-input-package}
Let $S$ be a quasi-compact quasi-separated spectral scheme. Then the
following data and properties are canonically available.
\begin{enumerate}[label=(\roman*)]
    \item The $\infty$-category
    \[
        \operatorname{Sm}^{\mathrm{fp}}_{/S}
    \]
    is essentially small and admits finite products and finite coproducts.

    \item Spectral Nisnevich distinguished squares generate a
    subcanonical Grothendieck topology
    \[
        \tau_{\mathrm{Nis}}
    \]
    on $\operatorname{Sm}^{\mathrm{fp}}_{/S}$.

    \item A presheaf of spaces satisfies Nisnevich descent if and only if
    it is reduced and carries every distinguished square to a Cartesian
    square.

    \item The brave new affine line
    \[
        \mathbb A^1_S
        =
        \Spec_S
        \operatorname{Sym}_{\mathcal O_S}(\mathcal O_S)
    \]
    is a smooth spectral $S$-scheme of finite presentation, equipped with
    zero and unit sections.

    \item The smooth site, Nisnevich topology, distinguished squares, and
    brave new affine line are compatible with base change between quasi-compact quasi-separated
    spectral bases; the underlying geometric constructions admit arbitrary base change.

    \item Open immersions, closed complements, \`etale morphisms, smooth
    morphisms, vector bundles, infinitesimal extensions, and regular closed
    immersions satisfy the base-change and deformation properties proved
    in the preceding chapters and in this chapter.
\end{enumerate}
Consequently the triple
\[
    \left(
        \operatorname{Sm}^{\mathrm{fp}}_{/S},
        \tau_{\mathrm{Nis}},
        \mathbb A^1_S
    \right)
\]
is defined functorially on qcqs spectral bases and supplies the geometric input for brave
new motivic homotopy theory.
\end{theorem}

\begin{proof}
Essential smallness and finite products and coproducts are
\Ref{prop:smooth-site-essentially-small},
\Ref{prop:smooth-site-finite-products}, and
\Ref{prop:smooth-site-finite-coproducts}. The Nisnevich topology and its
subcanonicity are
\Ref{def:spectral-nisnevich-topology} and
\Ref{thm:spectral-nisnevich-subcanonical}; the excision criterion is
\Ref{thm:nisnevich-excision-criterion}. Smoothness of the affine line is
\Ref{prop:affine-line-in-smooth-site}. Base-change compatibility is
\Ref{prop:smooth-site-geometric-base-change},
\Ref{prop:nisnevich-square-base-change}, and
\Ref{prop:affine-line-projection-base-change}. The remaining geometric
properties are the permanence and deformation theorems established in
\Ref{sec:closed-immersions-complements},
\Ref{sec:quasi-smooth-closed-immersions}, and
\Ref{sec:smooth-etale-lifting-neighbourhoods}.
\end{proof}

The constructions of the chapter may therefore be summarized by the
type-correct collection
\[
    \CAlg(\Sp)^{\mathrm{cn}}
    =
    \Alg_{\Comm^\otimes}(\Sp^\otimes)^{\mathrm{cn}},
    \qquad
    \Aff_{\Sp}
    :=
    \bigl(\CAlg(\Sp)^{\mathrm{cn}}\bigr)^{\op},
\]
\[
    \Aff_{\Sp}
    \longrightarrow
    \operatorname{Sch}^{\mathrm{sp}},
\]
\[
    \mathbf{Mod}_{\Sp}
    \colon
    \CAlg(\Sp)^{\mathrm{cn}}
    \longrightarrow
    \CAlg(\Pr^L_{\mathrm{st}}),
    \qquad
    A
    \longmapsto
    \Mod_A,
\]
\[
    \Spec_X
    \colon
    \operatorname{QCAlg}(X)^{\mathrm{cn},\op}
    \longrightarrow
    \operatorname{Sch}^{\mathrm{sp}}_{/X},
\]
and
\[
    \mathbb A^1_S
    \in
    \operatorname{Sm}^{\mathrm{fp}}_{/S},
    \qquad
    \bigl(
        \tau_{\mathrm{Nis}},
        W_{\mathbb A^1,S}
    \bigr).
\]
The first two displays record operadic algebra, the definition of affine
spectral geometry, and passage to global spectral schemes. The third
records affine module semantics, and the fourth records relative affine
geometry. The final display records the brave new interval object,
the spectral Nisnevich topology, and the affine-line projection class
which will be used to define the two successive localizations of
presheaves.

The next chapter begins with the presentable presheaf
\(\infty\)-category
\[
    \PSh
    \bigl(
        \operatorname{Sm}^{\mathrm{fp}}_{/S}
    \bigr)
\]
and constructs the \(\infty\)-category of brave new motivic spaces by
first imposing Nisnevich descent and then imposing
\(\mathbb A^1\)-invariance. The first localization is determined by the
Grothendieck topology
\[
    \tau_{\mathrm{Nis}},
\]
or equivalently by the Nisnevich excision conditions established in
this chapter. The second localization inverts the morphisms in
\[
    W_{\mathbb A^1,S}.
\]
Equivalently, it imposes the equivalences
\[
    F(X)
    \xrightarrow{\simeq}
    F(X\times_S\mathbb A^1_S)
\]
for every
\[
    X
    \in
    \operatorname{Sm}^{\mathrm{fp}}_{/S}.
\]
The subsequent constructions develop pointing and stabilization.
